\section{Problém \ensuremath{\SAT}}\label{sec:sat}

V~podsekci \ref{subsec:cnf} jsme viděli,~že převod obecné formule do CNF \uv{hrubou silou} je problematický. Nová formule se totiž může až exponenciálně prodloužit. V~podsekci \ref{subsec:cnf_tseitin} jsme si proto ukázali Tseitinovu transformaci,~která pomocí nových proměnných sestrojí stejně splnitelnou CNF efektivně. Zkusme si tedy původní problém zjednodušit. Předpokládejme nyní,~že formule na vstupu jsou již v~CNF. Tento problém je v~informatice velmi význačný a~označuje se zkratkou \ensuremath{\SAT} (z~angl. \emph{satisfiability},~neboli splnitelnost).

\problem{Splnitelnost logické formule v~CNF (\ensuremath{\SAT})}{Logická formule $\varphi$ v~CNF.}{1,~pokud je $\varphi$ splnitelná,~jinak 0.}

Nyní se zaměříme na samotnou splnitelnost vstupní formule. V~praxi se používají tzv. \emph{\ensuremath{\SAT} solvery},~které rozhodují,~zda je formule v~CNF splnitelná. Ukážeme si klasický algoritmus \emph{DPLL},~jehož název připomíná příjmení Martina Davise,~Hilaryho Putnama,~George Logemanna a~Donalda Lovelanda. Jeho základní podoba pochází z~počátku 60.~let 20.~století.

\subsection{Algoritmus DPLL}\label{subsec:dpll}

Tento algoritmus je založený na rekurzivním postupu. Na vstupu algoritmu je libovolná formule $\varphi$ v~CNF,~kterou algoritmus postupně vyhodnocuje. Využívá dvojici procedur -- \textbf{Unit propagation} a~\textbf{Pure literal elimination}.

\subsubsection*{Unit propagation}

Tato procedura hledá v~dané formuli $\varphi$ tzv. \emph{jednotkovou klauzuli}. To je klauzule,~která obsahuje \textbf{pouze jeden literál}. Například formule $\omega$ níže obsahuje jednotkovou klauzuli.
\[\omega(x,y)=(x\lor\neg y)\land \underbrace{\neg y}_{\text{jednotková kl.}}\land (\neg x\lor y).\]
Pokud se taková klauzule nachází ve formuli,~lze toho využít,~neboť pro danou proměnnou existuje právě jedno ohodnocení,~které tuto klauzuli splní. Obecně pokud formule $\varphi$ obsahuje jednotkovou klauzuli s~proměnnou $x$,~pak mohou nastat dva případy. Označme $\varphi^\prime$ zbytek klauzulí ve formuli $\varphi$.
\begin{enumerate}
    \item Pokud má formule tvar $\varphi=x\land\varphi^\prime$,~pak nutně $x=1$.
    \item Jinak pokud má formule tvar $\varphi=\neg x\land\varphi^\prime$,~pak musí být $x=0$.
\end{enumerate}
Fakt,~že jsme schopni hned rozhodnout,~jaké ohodnocení musí mít daná proměnná,~je v~tomto případě hodně ku prospěchu,~neboť formuli můžeme zjednodušit. Protože ohodnocení dané jednotkové klauzule již známe,~můžeme ji z~formule $\varphi$ odstranit. Zároveň můžeme \textbf{odstranit všechny klauzule},~které se díky proměnné $x$ vyhodnotí na $1$,~neboť již nemohou nijak ovlivnit logickou hodnotu zbylých klauzulí.

Např. pokud formule obsahuje literál $x$ v~jednotkové klauzuli,~pak hodnota proměnné musí být $1$. Tzn. klauzule obsahující $x$ se vyhodnotí též na $1$. tj. 
\[(\dots\lor x\lor\dots)=1.\]
Naopak pokud obsahuje literál $\neg x$,~pak nutně $x=0$ a~tedy
\[(\dots\lor \neg x\lor\dots)=1.\]
Dále ještě odstraníme výskyty literálu $x$,~resp. $\neg x$ v~\emph{opačné polaritě},~tedy negaci původního literálu. Tzn. pokud literál v~jednotkové klauzuli byl $x$,~pak opačná polarita literálu je $\neg x$ a~naopak. Důvod,~proč jej můžeme odstranit,~je ten,~že v~rámci kterékoliv jiné klauzule se daný literál v~opačné polaritě vyhodnotí vždy na $0$. Protože je však spojen s~ostatními literály pomocí $\lor$,~nemůže nijak ovlivnit výslednou logickou hodnotu klauzule a~tudíž jej můžeme odstranit. Tedy např. pokud $x=1$,~pak můžeme klauzule tvaru
\[(\dots\lor y \lor \underbrace{\neg x}_{\text{odstraníme}} \lor z~\lor \dots)\]
zjednodušit na
\[(\dots\lor y \lor z~\lor \dots).\]
Podobně pro $x=0$:
\[(\dots\lor y \lor x \lor z~\lor \dots)=(\dots\lor y \lor z~\lor \dots).\]

\begin{example}
    Máme formuli
    \[\chi(x,y,z)=(x\lor y)\land y\land (z\lor \neg y)\land(\neg x\lor \neg y\lor \neg z).\]
    Aplikujme na ni proceduru \emph{Unit propagation}. Můžeme si všimnout,~že klauzule $y$ je jednotková a~proměnná $y$ tedy musí být nastavena na $1$.

    Protože $y=1$,~je zároveň splněna i~klauzule $x\lor y$,~takže ji můžeme odstranit. Zároveň odstraníme literál $\neg y$ ze zbývajících klauzulí $z\lor\neg y$ a~$\neg x\lor\neg y\lor\neg z$.
    Celkově dostaneme novou zjednodušenou formuli
    \[\chi^\prime(x,y,z)=z\land(\neg x\lor \neg z).\]
\end{example}

\subsubsection*{Pure literal elimination}

Ve vstupní formuli $\varphi$ může nastat situace,~kdy se zde určitá proměnná nachází pouze v~jedné polaritě. Tzn. pokud se ve formuli nachází proměnná $x$,~pak se v~ní nikde nenachází $\neg x$ a~naopak pokud se v~ní nachází $\neg x$,~pak v~ní nikde není $x$. V~obou případech můžeme rovnou určit hodnotu dané proměnné,~neboť tím splníme všechny klauzule,~v~nichž se vyskytuje. Ty můžeme odstranit,~neboť jsou tím pádem splněny.
\[\dots\land\overbrace{(\dots\lor \underbrace{x}_{1}\lor\dots)}^{1}\land\dots\land\overbrace{(\dots\lor \underbrace{x}_{1}\lor\dots)}^{1}\land\dots\]
Podobně pro negaci,~tj.
\[\dots\land\overbrace{(\dots\lor \neg\underbrace{x}_{0}\lor\dots)}^{1}\land\dots\land\overbrace{(\dots\lor \neg\underbrace{x}_{0}\lor\dots)}^{1}\land\dots\]

Poté,~co jsme si vysvětlili základní dvojici procedur,~se nyní nabízí trojice případů,~které je ještě třeba prozkoumat.
\begin{itemize}
    \item \emph{Co když postupným odstraňováním proměnných mi ve formuli vznikne prázdná klauzule?} Taková situace může nastat v~případě procedury Unit propagation,~kdy odstraňujeme proměnné v~opačné polaritě. Pokud nastane situace,~že je nějaká klauzule prázdná,~pak to znamená,~že daná klauzule pod zvoleným ohodnocením \textbf{není splnitelná} a~tudíž ani celá formule. V~takovou chvíli můžeme prohlásit formuli za \textbf{nesplnitelnou pod daným ohodnocením}\footnote{Je potřeba si uvědomit,~že ohodnocení proměnných si během výpočtu volíme,~tedy v~takovém případě to ještě nutně neznamená,~že je formule celkově nesplnitelná. Pouze víme,~že dané ohodnocení určitě není splňující}.
    \item \emph{Co když odstraníme postupně všechny klauzule?} K~takové situaci může dojít,~pokud jsou všechny klauzule splněny. V~takovou chvíli určitě víme,~že je původní formule \textbf{splnitelná}.
    \item \emph{Co když po aplikaci obou procedur stále nevíme,~zda je $\varphi$ splnitelná?} V~takovou chvíli nám nezbývá nic jiného než si zvolit nějakou dosud neohodnocenou proměnnou a~prozkoumat obě možnosti jejího ohodnocení. Vybereme tedy proměnnou $x$ a~zkusíme ji nastavit jednou na hodnotu $1$ a~podruhé na hodnotu $0$. Tím můžeme formuli $\varphi$ zjednodušit a~na nově vzniklou formuli $\varphi^\prime$ rekurzivně aplikovat stejný algoritmus. Pokud lze splnit zbytek formule $\varphi^\prime$,~pak lze splnit i~původní formuli $\varphi$.
    
    Toho můžeme dosáhnout tak,~že k~formuli $\varphi$ přidáme uměle jednotkovou klauzuli $x$,~resp. v~druhém případě $\neg x$. Tuto jednotkovou klauzuli si pak vybere procedura Unit propagation,~která proměnnou ohodnotí. Zavoláme tedy algoritmus DPLL rekurzivně na formule
    \[\varphi\land x\quad\text{a}\quad\varphi\land\neg x.\]
\end{itemize}

S~těmito informacemi můžeme dát dohromady pseudokód algoritmu.
\begin{algorithm}
    \caption{DPLL}
    \label{alg:dpll}
    \KwIn{Formule $\varphi$ v~CNF}
    \While{\textup{existuje jednotková klauzule $(\ell)$ ve formuli $\varphi$}}{
        $\varphi\gets\text{UnitPropagation}(\ell,~\varphi)$
    }
    \While{\textup{existuje literál $\ell$ v~jediné polaritě ve $\varphi$}}{
        $\varphi\gets\text{PureLiteralElimination}(\ell,\varphi)$
    }
    \tcp{Všechny klauzule jsou pravdivé}
    \lIf{je $\varphi$ prázdná}{\Return{$1$}}

    \tcp{Všechny literály v~klauzuli jsou nepravdivé}
    \lIf{\textup{obsahuje $\varphi$ prázdnou klauzuli}}{\Return{$0$}}
    \tcp{Zkusíme obě ohodnocení proměnné $x$}
    Vyber dosud neohodnocenou proměnnou $x$\\

    \Return{$\textup{DPLL}(\varphi\land x)\lor\textup{DPLL}(\varphi\land\neg x)$}

    \KwOut{$1$,~je-li formule $\varphi$ splnitelná,~jinak $0$.}
\end{algorithm}

Zkusme si algoritmus odsimulovat na příkladu.
\begin{example}
    Máme formuli
    \[\gamma(x_1,x_2,x_3,x_4)=(x_1\lor x_3)\land\neg x_4\land(\neg x_2\lor\neg x_3)\land x_1\land(\neg x_2\lor x_3)\land(\neg x_1\lor\neg x_3\lor x_3).\]
    Chceme pomocí algoritmu DPLL rozhodnout,~zda je splnitelná.

    Když se podíváme na pseudokód \ref{alg:dpll} výše,~tak první máme aplikovat proceduru \emph{Unit propagation}. Hledáme tedy jednotkové klauzule ve formuli $\gamma$,~dokud nějaké existují.
    \begin{itemize}
        \item První jednotková klauzule,~které si můžeme všimnout,~je $\neg x_4$. Hodnota proměnné $x_4$ tedy nutně musí být $0$. Klauzuli tedy odstraníme. Proměnná se nachází ve formuli pouze jednou,~tedy daný literál v~opačné polaritě nikde odstraňovat nemusíme. Máme nyní formuli
        \[\gamma^\prime(x_1,x_2,x_3,x_4)=(x_1\lor x_3)\land(\neg x_2\lor\neg x_3)\land x_1\land(\neg x_2\lor x_3)\land(\neg x_1\lor\neg x_3\lor x_3).\]
        \item Další jednotková klauzule je $x_1$. Proměnná $x_1$ musí být nutně $1$ a~klauzuli můžeme odstranit. Dále si však můžeme všimnout,~že tím pádem splněna i~klauzule $(x_1\lor x_3)$. Tu můžeme též odstranit. Literál $x_1$ se dále nachází v~opačné polaritě v~klauzuli $(\neg x_1\lor\neg x_3\lor x_3)$. Literál $\neg x_1$ můžeme také odstranit z~formule. Tím se nám formule $\gamma^\prime$ zjednoduší na
        \[\gamma^{\prime\prime}(x_1,x_2,x_3,x_4)=(\neg x_2\lor\neg x_3)\land (\neg x_2\lor x_3)\land(\neg x_3\lor x_3).\]
    \end{itemize}
    Nově vzniklá formule $\gamma^{\prime\prime}$ již žádné další jednotkové klauzule neobsahuje. Dále tedy můžeme aplikovat proceduru \emph{Pure literal elimination}.
    \begin{itemize}
        \item Formule $\gamma^{\prime\prime}$ obsahuje pouze jeden literál v~jediné polaritě,~a~to $\neg x_2$. Tím pádem hodnotu $x_2$ můžeme nastavit na $0$,~čímž splníme dvojici klauzulí $(\neg x_2\lor\neg x_3)$ a~$(\neg x_2\lor x_3)$. Ty můžeme odstranit,~čímž dostaneme formuli
        \[\gamma^{\prime\prime\prime}(x_1,x_2,x_3,x_4)=\neg x_3\lor x_3.\]
    \end{itemize}
    Protože formule $\gamma^{\prime\prime\prime}$ není ani prázdná,~ani neobsahuje prázdnou klauzuli,~vybereme náhodně jeden literál a~prozkoumáme jeho obě možná ohodnocení. V~tomto případě nám zbývají pouze literály $x_3$ a~$\neg x_3$. Vyberme např. literál $\neg x_3$. Nyní máme sestrojit dvojici nových formulí $(\neg x_3\lor x_3)\land\neg x_3$ a~$(\neg x_3\lor x_3)\land x_3$,~na které rekurzivně zavoláme DPLL.
    \begin{itemize}
        \item \textbf{Formule $(\neg x_3\lor x_3)\land\neg x_3$.} Opět začneme aplikací Unit propagation. K~původní formuli $\gamma^{\prime\prime\prime}$ jsme uměle přidali jednotkovou klauzuli $\neg x_3$. Z~ní je zjevné,~že $x_3$ musí být $0$. Tím zároveň splníme i~klauzuli $(\neg x_3\lor x_3)$,~kterou můžeme odstranit,~čímž dostáváme prázdnou formuli. Pure literal elimination aplikovat nemůžeme. Protože je nově vzniklá formule prázdná,~vrátíme na výstup $1$.
        \item \textbf{Formule $(\neg x_3\lor x_3)\land x_3$.} I~zde začneme eliminací jednotkových klauzulí. Tu zde představuje literál $x_3$,~jehož ohodnocení tedy musí být $1$. Tím opět splníme i~klauzuli $(\neg x_3\lor x_3)$,~jejímž odstraněním dostaneme zase prázdnou formuli. Tedy opět vrátíme $1$.
    \end{itemize}
    Obě ohodnocení pro $x_3$ jsou tedy splňující a~tedy i~původní formule $\gamma$ \emph{je splnitelná}.
\end{example}

Ukažme si ještě jeden jednodušší příklad.
\begin{example}
    Mějme formuli $\tau(x,y,z)=x\land (\neg x\lor y\lor z)\land \neg x$. Chceme opět pomocí DPLL určit splnitelnost.

    Jako první jednotkovou klauzuli můžeme vzít $x$. Tedy můžeme ji odstranit společně se všemi výskyty daného literálu v~opačné polaritě,~tj. $\neg x$. Tzn. dostaneme formuli
    \[\tau^\prime(x,y,z)=(y\lor z)\land().\]
    Všimněte si,~že odstraněním literálu $\neg x$ jsme dostali ve formuli prázdnou klauzuli\footnote{Z~jednotkové klauzule jsme odebrali její jediný literál v~rámci procedury Unit propagation; klauzule samotná ve formuli zůstala.} $()$. Aplikací Pure literal elimination se zbavíme i~literálů $y$ a~$z$.
    \[\tau^{\prime\prime}(x,y,z)=()\]
    Formule $\tau^{\prime\prime}$ obsahuje pouze prázdnou klauzuli. V~tuto chvíli algoritmus prohlásí formuli $\tau$ za \emph{nesplnitelnou}.
\end{example}

Poslední otázka,~která se nabízí,~je (jako obvykle) časová složitost algoritmu. Vezmeme-li v~potaz nejhorší případ,~časová složitost bohužel nevychází úplně příznivě.
\begin{theorem}[Složitost DPLL]
    Nechť formule $\varphi$ obsahuje $v$ proměnných a~její délka je $s$. Jednoduchá implementace DPLL doběhne v~čase $\bigO(s2^v)$. Při kopírování formule v~každé úrovni rekurze spotřebuje nejvýše $\bigO(sv)$ pomocné paměti.
\end{theorem}
\begin{proof}
    V~nejhorším případě nepomůže ani jednotková propagace,~ani eliminace čistých literálů. Algoritmus pak zvolí proměnnou a~vytvoří dvě větve,~jednu pro hodnotu $0$ a~druhou pro hodnotu $1$. Hloubka stromu je nejvýše $v$ a~počet jeho uzlů nejvýše
    \[1+2+4+\dots+2^v=2^{v+1}-1=\bigO(2^v).\]
    V~každém uzlu lze naivně projít celou formuli v~čase $\bigO(s)$,~odkud plyne čas $\bigO(s2^v)$. Rekurze probíhá do hloubky,~takže současně uchováváme nejvýše $v$ kopií formule délky nejvýše $s$; to dává $\bigO(sv)$ pomocné paměti.
\end{proof}
\notebox{Praktické \ensuremath{\SAT} solvery formuli při každém volání celou nekopírují. Změny zapisují do pomocných struktur a~při návratu je vracejí zpět,~takže paměťový odhad lze podstatně zlepšit.}

S~exponenciální časovou složitostí se sotva spokojíme,~avšak v~tomto případě je dobré si říci,~že zkoumáme složitost v~nejhorším případě a~v~praxi bývají zkoumané logické formule poměrně rychle řešitelné pomocí DPLL.
\importantbox{Problém je však dodnes ten,~že na rozhodnutí splnitelnosti logické formule (i~těch v~CNF) není známý žádný algoritmus,~který by běžel v~polynomiálním čase,~tj. $\bigO(n^k)$,~pro nějaké $k$. K~tomuto tématu se ještě vrátíme v~sekci \ref{sec:p_np}.}
